\section{FISH}
\label{sec:method}
FISH currently targets ROS~2 Humble and has three parts: \emph{instrumentation} of the ROS~2 client libraries and the GPU run-time, \emph{acquisition}, which runs an unmodified application under it, and \emph{extraction}, which turns the trace into the model of Fig.~\ref{fig:model}. The first challenge of Section~\ref{sec:goal} is met by what is recorded at initialisation and by how acquisition makes sure nothing escapes it; the second by what is recorded at execution and by how extraction ties the records together.

\noindent\textbf{The architecture.} FISH builds on \texttt{ros2\_tracing}~\cite{bedard-2022-ros2tracing}, whose initialisation events record each object as it is created and whose execution events record what then happens to it, linked by the object's handle. To the 14 stock events of Humble, FISH adds 28. From initialisation come the four things of Section~\ref{sec:goal}: introspection events record each executor's type and thread count and each callback's group; finalize events close the lifetimes the stock events open, so a reused handle is not mistaken for its predecessor; a timer event binds a timer to its node and period; registration events give rclpy nodes the same chain as rclcpp; and the intra-process subscription waitable makes deliveries that never reach \texttt{rcl} visible, completing the five kinds the executor dispatches. Link events at the publish, receive and request sites tie each message to the callback that issued it. For nothing to escape, the instrumented libraries replace the stock ones on the host or in a container, and the FISH entry point wraps \texttt{ros2 launch}. Because CUPTI sees only processes started under the profiler, the entry point reads the launch description and starts the containers that will host GPU components under \texttt{nsys} from the outset; standalone GPU processes are caught by the FISH daemon, which polls \texttt{/proc} and restarts any process holding a CUDA context under the profiler. Once the node list has been unchanged over five polls and every component is loaded, the daemon starts the input. Initialisation events then create the vertices of Fig.~\ref{fig:model}a and the message and service edges the wiring declares, and the result is checked against what the traced binaries create: a node constructor must yield one node, seven entities and seven callbacks, and the architecture is complete only if every such expectation appears in it.

\noindent\textbf{The timelines.} Start and end events for every kind of callback record each execution with the thread it ran on, and the scheduler's context switches split the time between them into on core, preempted and blocked. GPU work is recorded by CUPTI through \texttt{nsys}, aligned to the LTTng clock at ingest and bound through correlation ids to the API call, and through the calling thread's active callback to the callback that issued it. Phase boundaries come from the trace: initialisation ends when every periodic timer has fired once, warm-up when every recurring callback has completed its first execution, and a callback that never ran is kept as such. Extraction then finds the callback that actually issued each declared edge, adds the state and join edges that only running reveals (Fig.~\ref{fig:model}c), and attaches the measurements. Below the callback, one that submits GPU work owns a typed DAG per invocation (Fig.~\ref{fig:model}d): CPU segments, kernels and copies, ordered by launch, CPU order and stream FIFO. Invocations with the same node sequence form one shape, and folding the steady-state shapes onto their longest common subsequence, each branch carrying the share of invocations that take it, gives the callback's conditional graph.

\noindent\textbf{One model.} Every edge is lifted along the containment, so a message between two callbacks is also one between their nodes and their executors (Fig.~\ref{fig:model}b); the expansion operator is the inverse step, replacing a vertex by its children and its edges by theirs, which is how the graph is read at mixed granularity, one executor opened into its nodes while the rest stay whole, and how a chain is followed across executors, nodes and processes and into the GPU work of a callback. FISH-eye, a browser-based viewer, renders the model at any granularity the operator allows (Fig.~\ref{fig:screens}).
